\section{第9章の解答}\label{解答:第9章}

\paragraph{9.1}
まず交点と上下関係を調べる.\,$x^2=3x$ より交点の $x$ 座標は $0,3$ であり,\,$0\le x\le3$ では $3x-x^2=x(3-x)\ge0$ である.したがって面積は,上の曲線から下の曲線を引いて
\[
\begin{aligned}
A&=\int_0^3(3x-x^2)\dd x\\
 &=\left[\frac{3x^2}{2}-\frac{x^3}{3}\right]_0^3\\
 &=\frac{27}{2}-9\\
 &=\frac92.
\end{aligned}
\]
この区間では両曲線とも $x$ 軸の上側にある.回転後の断面は,外半径 $3x$,内半径 $x^2$ の円環なので,体積は
\[
\begin{aligned}
V&=\pi\int_0^3\{(3x)^2-(x^2)^2\}\dd x\\
 &=\pi\left[3x^3-\frac{x^5}{5}\right]_0^3\\
 &=\pi\left(81-\frac{243}{5}\right)\\
 &=\frac{162\pi}{5}.
\end{aligned}
\]
面積では高さの差を使い,体積では円板の面積の差を使う.帯の高さ $3x-x^2$ を1つの半径として2乗すると,別の断面を数えることになる.

\paragraph{9.2}
横の帯を使うので,交点も $y$ 座標で求める.
\[
y^2=3-2y
\quad\Longleftrightarrow\quad
(y+3)(y-1)=0.
\]
よって積分区間は $-3\le y\le1$ である.この間では
\[
(3-2y)-y^2=(1-y)(y+3)\ge0
\]
だから,直線 $x=3-2y$ が右側,放物線 $x=y^2$ が左側にある.求める面積は
\[
\begin{aligned}
A&=\int_{-3}^1(3-2y-y^2)\dd y\\
 &=\left[3y-y^2-\frac{y^3}{3}\right]_{-3}^1\\
 &=\left(3-1-\frac13\right)-(-9-9+9)\\
 &=\frac{32}{3}.
\end{aligned}
\]
横の帯の幅を右端から左端を引いた差で表したため,放物線を2つの平方根に分ける必要がない.

\paragraph{9.3}
各 $x$ で線分 $-1\le y\le x$ は回転軸を含むので,回転後の断面に穴はない.半径は,上端までの距離 $x$ と下端までの距離 $1$ のうち大きいほうである.したがって外半径は
\[
R(x)=
\begin{cases}
1,&0\le x\le1,\\
x,&1\le x\le2.
\end{cases}
\]
外半径を決める端が交代するのは $x=1$ である.断面の円板を積み重ねると
\[
\begin{aligned}
V&=\pi\int_0^1 1\dd x+\pi\int_1^2 x^2\dd x\\
 &=\pi+\pi\left[\frac{x^3}{3}\right]_1^2\\
 &=\pi+\frac{7\pi}{3}\\
 &=\frac{10\pi}{3}.
\end{aligned}
\]
上側と下側を別々に回した円板は重なる.2つの円板の面積を加えるのでなく,大きいほうの円板を一度だけ数えている.

\paragraph{9.4}
円筒殻を使うときは,\,$y$ 軸に平行な縦の帯を取る.回転半径は $x$,層の高さは $1-x^2$ だから,
\[
\begin{aligned}
V&=2\pi\int_0^1x(1-x^2)\dd x\\
 &=2\pi\left[\frac{x^2}{2}-\frac{x^4}{4}\right]_0^1\\
 &=2\pi\left(\frac12-\frac14\right)\\
 &=\frac{\pi}{2}.
\end{aligned}
\]
円板を使うときは,\,$y$ 軸に垂直な横の帯を取る.\,$0\le y\le1$ で,元の領域の左右の端は $x=0$ と $x=\sqrt y$ である.回転後の断面は半径 $\sqrt y$ の円板なので,
\[
V=\pi\int_0^1(\sqrt y)^2\dd y
 =\pi\left[\frac{y^2}{2}\right]_0^1=\frac{\pi}{2}.
\]
2つの式では積分変数と分け方が異なるが,どちらも同じ立体を重複なく数えている.

\paragraph{9.5}
グラフの弧長では $\sqrt{1+y'^2}$ を積分する.まず合成関数の微分により
\[
y'=\frac13\cdot\frac32(x^2+2)^{\frac12}\cdot2x
 =x\sqrt{x^2+2}.
\]
このまま根号を積分しようとせず,中身を整理する.
\[
1+y'^2=1+x^2(x^2+2)=x^4+2x^2+1=(x^2+1)^2.
\]
$x^2+1>0$ なので,平方根は $|x^2+1|=x^2+1$ となる.したがって
\[
L=\int_0^1(x^2+1)\dd x
 =\left[\frac{x^3}{3}+x\right]_0^1=\frac43.
\]
弧長の根号が完全平方になっていることが,計算を簡単にする理由である.

\paragraph{9.6}
速度を展開すると $v(t)=3t^2-9t+6$ である.これを積分し,\,$s(0)=0$ を使えば位置は
\[
s(t)=t^3-\frac92t^2+6t
\]
となる.終点の位置は
\[
s(3)=27-\frac{81}{2}+18=\frac92.
\]
道のりを求めるには,速度の符号を区別する.\,$v(t)=3(t-1)(t-2)$ は,\,$0\le t<1$ と $2<t\le3$ で正,\,$1<t<2$ で負である.折り返す位置は
\[
s(1)=1-\frac92+6=\frac52,\qquad
s(2)=8-18+12=2.
\]
各区間では一方向に動くので,道のりはそれぞれの位置の差の絶対値の和になり,
\[
D=\left|\frac52-0\right|
 +\left|2-\frac52\right|
 +\left|\frac92-2\right|
 =\frac52+\frac12+\frac52=\frac{11}{2}.
\]
これは $\int_0^1v\dd t-\int_1^2v\dd t+\int_2^3v\dd t$ と同じである.全時間は $3$ なので,平均速度と平均の速さはそれぞれ
\[
\frac{s(3)-s(0)}{3}=\frac32,\qquad
\frac{D}{3}=\frac{11}{6}
\]
となる.戻った分は変位では差し引かれ,道のりでは加えられるため,この運動では平均の速さのほうが大きい.

\paragraph{9.7}
関数の平均値は積分を区間の長さで割って求める.
\[
\overline f=\frac12\int_0^2(1+x^2)\dd x
 =\frac12\left[x+\frac{x^3}{3}\right]_0^2
 =\frac12\left(2+\frac83\right)=\frac73.
\]
$f(c)=\overline f$ となる条件は
\[
1+c^2=\frac73
\quad\Longleftrightarrow\quad
c^2=\frac43.
\]
$c\in[0,2]$ だから負の解は除き,
\[
c=\frac{2}{\sqrt3}=\frac{2\sqrt3}{3}
\]
を得る.区間の中央 $c=1$ では $f(1)=2$ なので,平均値を取る点とは異なる.

\paragraph{9.8}
三角形は $0\le x\le2$,\,$0\le y\le\frac{x}{2}$ と表せる.まず面積を求めると
\[
A=\int_0^2\frac{x}{2}\dd x
 =\left[\frac{x^2}{4}\right]_0^2=1.
\]
縦の帯の中心の高さは $\frac{x}{4}$ であり,帯の面積はほぼ $\frac{x}{2}\dd x$ だから,面積重心の高さは
\[
\begin{aligned}
\overline y_{\mathrm{面}}
 &=\frac1A\int_0^2\frac{x}{4}\cdot\frac{x}{2}\dd x\\
 &=\frac18\left[\frac{x^3}{3}\right]_0^2\\
 &=\frac13.
\end{aligned}
\]
三角形は軸の上側にあり,軸は境界の1辺に重なるが内部を横切らない.したがって体積の定理が使えて,
\[
V=2\pi\overline y_{\mathrm{面}}A
 =2\pi\cdot\frac13\cdot1=\frac{2\pi}{3}.
\]

全表面積では,境界の3本の線分を別々に考え,それぞれの線重心を使う.一様な線分の線重心は中点なので,次のようになる.

$x$ 軸上の線分は長さ $2$,線重心の高さ $0$ である.回転しても軸上の線分のままで,面積のある面を作らない.その寄与は $2\pi\cdot0\cdot2=0$ である.

縦の線分 $x=2$,\,$0\le y\le1$ は長さ $1$,線重心の高さ $\frac12$ である.回転すると半径 $1$ の円板になり,面の定理から
\[
S_{\mathrm{底面}}=2\pi\cdot\frac12\cdot1=\pi
\]
となる.端点が回転軸に触れていても,面を重複して作ることはない.

斜めの線分 $y=\frac{x}{2}$,\,$0\le x\le2$ は長さ $\sqrt5$,線重心の高さ $\frac12$ である.線重心を積分で確かめると,\,$\dd s=\frac{\sqrt5}{2}\dd x$ より
\[
\overline y_{\mathrm{線}}
 =\frac1{\sqrt5}\int_0^2\frac{x}{2}\cdot\frac{\sqrt5}{2}\dd x
 =\frac14\left[\frac{x^2}{2}\right]_0^2=\frac12.
\]
この線分が作る円すいの側面の面積は
\[
S_{\mathrm{側面}}=2\pi\cdot\frac12\cdot\sqrt5=\pi\sqrt5.
\]
底面と側面は円周で接するが,面積のある部分では重ならない.よって全表面積は
\[
S=\pi+\pi\sqrt5=\pi(1+\sqrt5)
\]
である.体積に使った面積重心の高さ $\frac13$ と,境界の各線分に使った線重心の高さを混同しない.
